% =========================================================================
%  Ejercicio 2b - Ensayo INDIVIDUAL sobre el articulo
%  Responsable: Diego Hernández Gómez
%
%  NOTA: Este ensayo es INDIVIDUAL, uno por cada integrante del equipo.
% =========================================================================

\subsection*{Ensayo individual — Diego Hernández Gómez}

\vspace{0.3cm}

\noindent
\textcolor{unamblue}{\rule{\textwidth}{1.2pt}}

\vspace{0.3cm}

\subsubsection*{\textcolor{unamblue}{Introducción}}

El artículo habla de cómo estamos llegando a una nueva revolución industrial donde nosotros somos el producto o mejor dicho los datos que generamos, habla de cómo estos datos han ido ganando cada vez más importancia para las empresas y esto es bastante interesante e importante de conocer, la mayoría de las personas o no conoce de esto o no le impronta, usualmente los usuarios ceden sus datos a empresas que tienen servicios en internet como Youtube,Google,etc; mediante sus términos y condiciones. A mi parecer no lo veo mal, sería como el costo de los servicios, pero concuerdo con el autor de que debemos de asegurarnos que no se haga mal uso de esos datos.

\vspace{0.1cm}

\subsubsection*{\textcolor{unamblue}{Desarrollo}}

La autora principalmente argumenta que los datos personales funcionan como el nuevo activo monetario que sustenta la economía digital bajo una falsa apariencia de gratuidad a cambio de la privacidad. Asimismo, se señala una profunda asimetría de poder frente a las corporaciones tecnológicas y un incremento de la vigilancia mediante Internet of Things, lo que exige un empoderamiento del usuario y regulaciones como el GDPR. 

\vspace{0.2cm}

Lo importante de esto se empezar a conocer lo que se hace con nuestra información, el desequilibrio extremo entre empresas y usuarios,las empresas saben todo sobre nosotros, pero nosotros no sabemos qué hacen con nuestros datos, un ejemplo de esto podrían ser los algoritmos de redes sociales, estos aprenden de nuestros gustos opiniones ideas, algunos funcionan claramente para hacernos interactuar al recomendar o directamente mostrarnos contenido que nos guste o nos desagrade, pero otros usan los datos para censurar lo que las empresas no quieran que se vea; otro ejemplo serían las IAs, una vez comenzó su auge muchos de los datos de las personas en redes sociales u otras plataformas, fueron usados para entrenarlas, un ejemplo claro fueron las IAs para generar imágenes o videos, se usaron fotos, dibujo y videos subidos por usuarios sin su consentimiento y no pudieron hacer nada para impedirlo. Así hay otros ejemplos de cómo las empresas usan los datos de los usuarios en su contra o sin su permiso.

\vspace{0.2cm}

La autora principalmente buscaba que nos empezáramos a preguntar cuántos de estos datos hemos dado sin darnos cuenta y si siquiera sabemos lo que hacen con ellos, además de darnos cuenta cuánto realmente vale nuestros datos para las empresas. Esto se relaciona bastante con Fundamentos de Bases de Datos, pues esto se técnicamente lo que buscamos aprender y entender, el cómo se maneja y filtra de esa cantidad de datos demencial, aquellos datos que nos son útiles a nosotros o a las empresas para nuestros propósitos, pero también nos hace ver el valor de la información con la que tendríamos que trabajar si nos centramos en algo como esto.

\vspace{0.2cm}

\subsubsection*{\textcolor{unamblue}{Conclusión}}

Si bien no veo mal la recolección de datos por las empresas mal, estoy de acuerdo que el saber y tener formas de controlar qué datos compartimos y  lo que se hace con ello. Lo que me llevo de este artículo es a valorar más mis datos y los de los demás; y a prestar más atención a cómo está actualmente y en un futuro la recolección y uso de los datos generados por las personas.

\vspace{0.2cm}

\noindent
\textcolor{unamblue}{\rule{\textwidth}{1.2pt}}
